\section{Convergence of the companion remainder}\label{sec:companion}
We use $\delta=t-3\in\Omega=\C\setminus(-\infty,-3]$.
All analytic logarithms in this section vanish at $\delta=0$.
Distinguish the original real-line weight $w_y$ in
Section~\ref{sec:boundary} from the angular multiplier
\begin{equation}\label{wang}
 w(\theta)=\frac1{1+2\cosh(\xi(\theta)/3)},\quad
 \xi(\theta)=2\log\cot(\theta/2),\quad s_0(\theta)=\sqrt{w(\theta)}.
\end{equation}
Thus $0\le w\le1/3$ and $w=O(d^{2/3})$ for
$d(\theta)=\min(\theta,\pi-\theta)$.
Let $U:\ell^2(\N_0)\to L^2(0,\pi)$ synthesize the real normalized
Jacobi angular basis with parameters $(4/3,2/3)$, and let $H$ be the
orthogonal Hardy projection obtained from the negative real half-line
under Fourier and angular transformation. Explicitly it is unitarily
conjugate to the Fourier image of multiplication by $\one_{(-\infty,0)}$.
Write $P_n$ for the first $n$ coordinate projection and $J_n=UP_nU^*$.
Set
\[
 X=M_{s_0}HU,\quad Y=M_{s_0}(I-H)U,\quad
 A_n=P_nX^*XP_n,\quad B_n=P_nY^*YP_n.
\]
Here $A_n,B_n$ act on $\Ran P_n$. Define
\[
 Q_n(\delta)=\det(I+\delta A_n)=\det(I+\delta B_n),\quad
 S_n(\delta)=\det(I+\delta(A_n+B_n)),
\]
\[
 T_n(\delta)=\det_{\Ran J_n}(I+\delta J_nM_wJ_n).
\]
Equality of the two single determinants follows from complex
conjugation in the real Jacobi basis, which interchanges $H$ and $I-H$.
The supplied earlier proof gives, locally uniformly on $\Omega$,
\begin{equation}\label{inheritedfactors}
 \frac{D_n(t)/D_n(3)}{Q_n(t-3)}\longrightarrow\mathcal A(t),
 \qquad \frac{S_n(t-3)}{T_n(t-3)}\longrightarrow\Xi(t).
\end{equation}
For clarity, their complete operator definitions are
\begin{align}
 \mathcal A(t)&=\det{}_2\!\left(I+\delta(I+\delta C_0)^{-1}\mathcal E\right)
 \exp\!\left\{\delta\tr[(I+\delta C_0)^{-1}\mathcal S]\right\},
 \label{Arel}\\
 \Xi(t)&=\det\!\left[(I+\delta\mathcal D)(I+\delta M_w)^{-1}\right],
 \quad \mathcal D=HM_wH+(I-H)M_w(I-H).
 \label{Xidef}
\end{align}
In \eqref{Arel} the real-line operators are exactly those defined in
\eqref{kernelreal} below. The difference $\mathcal D-M_w$ is trace
class; the two factors are invertible on $\Omega$. Both $\mathcal A$
and $\Xi$ are nowhere zero, normalized to one at calibration, and
positive on $t>0$. The gamma-phase comparison has limiting relative
determinant one, which is already included in the first limit in
\eqref{inheritedfactors}.

\subsection{A moving-edge mixed determinant limit}
The weighted-wave result of \cite[Section 11]{r124} gives bounded maps
\[
 X_0e_j=\frac{s_0(\theta)e^{-i(j+3/2)\theta+ic_0}}{\sqrt{2\pi}},
 \qquad Y_0e_j=\frac{s_0(\theta)e^{i(j+3/2)\theta-ic_0}}{\sqrt{2\pi}},
 \quad c_0=11\pi/12,
\]
with $X-X_0,Y-Y_0\in\Stwo$. Importantly, $XY^*=0$ exactly,
since $UU^*=I$ and $H(I-H)=0$.
Define $M_n$ to be multiplication by $e^{-in\theta}$ and define the
Hilbert--Schmidt operator $K$ by its kernel
\begin{equation}\label{Kedge}
 K(\theta,\phi)=\frac{s_0(\theta)s_0(\phi)}{2\pi}
 \frac{e^{2ic_0-(3i/2)(\theta+\phi)}}{1-e^{-i(\theta+\phi)}}.
\end{equation}
At either singular corner its square is bounded by
$C\theta^{2/3}\phi^{2/3}/(\theta+\phi)^2$ or its reflection,
which is integrable after setting $\theta=ru$, $\phi=r(1-u)$.
Thus $K\in\Stwo$.

The geometric sum gives the exact identity
\[
 X_0P_nY_0^*=K-M_nKM_n.
\]
Since $M_n\rightharpoonup0$ by the Riemann--Lebesgue lemma and $K$
is compact, $M_nKM_n\to0$ strongly. Taking the strong limit gives
$X_0Y_0^*=K$. Set $E_X=X-X_0$, $E_Y=Y-Y_0$. Finite-rank approximation
in $\Stwo$ gives
\[
 XP_nY^*-X_0P_nY_0^*=E_XP_nY^*+X_0P_nE_Y^*
       \longrightarrow E_XY^*+X_0E_Y^*=-K
\]
in Hilbert--Schmidt norm. Consequently
\begin{equation}\label{movingK}
 K_n:=XP_nY^*=-M_nKM_n+o_{\Stwo}(1).
\end{equation}
In particular $\tr(A_nB_n)=\|K_n\|_2^2\to\|K\|_2^2$.
The moving term in \eqref{movingK} is retained; it does not tend to
zero in Hilbert--Schmidt norm.

Let $E_I$ extend functions on $(0,\pi)$ by zero to the circle
$(0,2\pi)$. Let $P_+^\circ$ project onto circle modes $j\ge1$ and
$P_-^\circ$ onto $j\le-1$, and set $G_\pm=M_{e^{\pm3i\theta/2}}$.
Define
\begin{equation}\label{edgeAB}
 A_e=M_{s_0}G_-E_I^*P_+^\circ E_IG_+M_{s_0},\qquad
 B_e=M_{s_0}G_+E_I^*P_-^\circ E_IG_-M_{s_0}.
\end{equation}
The strict mode cutoffs follow from $n-j=1,\ldots,n$ in the finite
wave sums, and yield
\begin{equation}\label{diagonallimits}
 M_n^*XP_nX^*M_n\longrightarrow A_e,\qquad
 M_nYP_nY^*M_n^*\longrightarrow B_e
\end{equation}
strongly. Here is why the fixed Hilbert--Schmidt errors do not affect
these limits. For fixed $h$, $P_nX^*M_nh$ is bounded and weakly null:
pairing with $g$ reduces to pairing $M_nh$ with $XP_ng\to Xg$ in
norm. Compactness of $E_X$ then gives
$E_XP_nX^*M_nh\to0$, and compactness also gives $E_X^*M_nh\to0$.
Expand $XP_nX^*-X_0P_nX_0^*=E_XP_nX^*+X_0P_nE_X^*$ to obtain
the first claim. The other is identical with reversed modulation.
Both limits are positive and bounded by $I/3$.

For $0\le T\le I/3$, $(I+\delta T)^{-1}$ is uniformly bounded
for $\delta$ in compact subsets of $\Omega$. Strong convergence
in \eqref{diagonallimits} gives strong convergence of these resolvents,
locally uniformly on each fixed vector. A finite parameter net and
the resolvent identity justify this uniformity.

Apply Sylvester's determinant identity to the block map with entries
$XP_n$ and $YP_n$, and take its upper-left Schur complement. With
\[
 R_{X,n}=(I+\delta XP_nX^*)^{-1},\qquad
 R_{Y,n}=(I+\delta YP_nY^*)^{-1},
\]
one obtains, with the order of factors as written,
\begin{equation}\label{schur}
 \frac{S_n(\delta)}{Q_n(\delta)^2}
 =\det[I-\delta^2R_{Y,n}K_n^*R_{X,n}K_n].
\end{equation}
Conjugate the product on the right by $M_n$, and use
\eqref{movingK}. Products of two Hilbert--Schmidt factors are trace
class, so all terms containing its $o_{\Stwo}(1)$ error tend to zero
in trace norm, uniformly on compact parameter sets. The main term is
\[
 \delta^2(M_nR_{Y,n}M_n^*)K^*(M_n^*R_{X,n}M_n)K.
\]
Strong resolvent convergence and finite-rank approximation of $K$
now give locally uniform trace-norm convergence. Hence
\begin{equation}\label{Psidef}
 \frac{S_n(\delta)}{Q_n(\delta)^2}\longrightarrow\Psi(\delta):=
 \det\!\left[I-\delta^2(I+\delta B_e)^{-1}K^*
                              (I+\delta A_e)^{-1}K\right].
\end{equation}
Every finite ratio is nonzero on $\Omega$, and $\Psi(0)=1$.
Hurwitz's theorem makes $\Psi$ nowhere zero. It is positive for
real $\delta>-3$, as the nonzero limit of positive ratios. Thus
\begin{equation}\label{mixedlimit}
 2\log Q_n(\delta)-\log S_n(\delta)\longrightarrow-\Log\Psi(\delta)
\end{equation}
locally uniformly, with the common normalization at zero. Alternatively,
the modulated positive block matrices have strong limit
$\bigl(\begin{smallmatrix}A_e&-K\\-K^*&B_e\end{smallmatrix}\bigr)$
bounded by $I/3$; its Schur complement directly proves invertibility
of the limiting factor in \eqref{Psidef}.

\subsection{The scalar constant: a polynomial theorem and its extension}
For an angular function $F(\theta)=g(\cos\theta)$ put
\[
 \mu(F)=\frac1\pi\int_0^\pi F(\theta)\,d\theta,\quad
 F_k=\frac1\pi\int_0^\pi F(\theta)\cos(k\theta)\,d\theta,\quad
 v(F)=\sum_{k\ge1}kF_k^2.
\]
For real $F$, the squares in $v$ are ordinary real squares.

\begin{proposition}\label{prop:scalarconstant}
If $F$ is real and angular H\"older of exponent
$1/2<\gamma\le1$, with $F(0)=F(\pi)=0$, then $v(F)<\infty$ and
\begin{equation}\label{scalarconstant}
 \log\det_{\Ran J_n}(J_ne^gJ_n)
     =(n+1)\mu(F)+\tfrac12v(F)+o(1).
\end{equation}
Here $g$ denotes multiplication by the function whose angular form is $F$.
\end{proposition}
\begin{proof}
The rank-$n$ Jacobi polynomial ensemble has density proportional to
its squared Vandermonde times the Jacobi weight. Andreief's finite
identity gives
\[
 \mathbb E e^{X_n(g)}=\det(J_ne^gJ_n),\quad
 \mathbb E X_n(g)=\tr(J_ngJ_n),\quad
 \Var_n(g)=\|(I-J_n)gJ_n\|_2^2.
\]
The earlier global kernel bound gives for every real angular
$C^\eta$ function $H$, $\eta>1/2$,
\begin{equation}\label{varianceholder}
 \Var_n(H)=\tfrac12\iint(H(\theta)-H(\phi))^2
                      |K_n(\theta,\phi)|^2\,d\theta\,d\phi
 \le C_\eta[H]_{C^\eta}^2,
\end{equation}
independently of $n$. Endpoint vanishing is unnecessary for this bound.

The polynomial input is Corollary 2.2 of Breuer--Duits
\cite{bd}: if the recurrence matrix has a Laurent right limit,
every fixed real polynomial statistic, centered by its exact mean,
converges to a centered Gaussian with variance
$\sum_{k\ge1}k\widehat p_k\widehat p_{-k}$, where the Fourier
coefficients are those of the polynomial applied to the Laurent symbol.
The fixed Jacobi recurrence here has off-diagonal limit $1/2$ and
diagonal limit zero, hence the unique full-sequence symbol
$(z+z^{-1})/2$. Its coefficients are bounded and the support is compact.
Thus this theorem applies to every real polynomial $p$, with variance
$v(p(\cos\theta))$. Only the polynomial theorem is used; its separate
$C^1$ extension is not applied to the fractional endpoint symbol.

A weak CLT alone is insufficient for the determinant. To obtain
exponential integrability, let $E_r=e^{rg/2}$ and
$T_r=J_ne^{rg}J_n$ on $\Ran J_n$. The isometry
$E_rJ_nT_r^{-1/2}$ has range projection
$Q_r=E_rJ_nT_r^{-1}J_nE_r$. For any bounded real multiplier $a$,
commutation with $E_r$ gives the exact identity
\[
 (I-Q_r)aE_rJ_nT_r^{-1/2}
   =(I-Q_r)E_r(I-J_n)aJ_nT_r^{-1/2}.
\]
For $r\ge0$, taking Hilbert--Schmidt norms gives
$\Var_{rg}(a)\le e^{r\operatorname{osc}(g)}\Var_n(a)$.
Taylor's integral formula for the logarithmic moment generating
function, also applied to $-g$, therefore yields for real $u$
\begin{equation}\label{exponentialbound}
 \log\mathbb E e^{u(X_n(g)-\mathbb EX_n(g))}
 \le \tfrac12u^2e^{|u|\operatorname{osc}(g)}\Var_n(g).
\end{equation}
For fixed polynomial $p$, the bound at $u=2$ and
\eqref{varianceholder} prove uniform integrability of
$e^{X_n(p)-\mathbb EX_n(p)}$. The polynomial CLT consequently implies
\begin{equation}\label{polynomialG}
 G_n(p):=\log\mathbb E e^{X_n(p)-\mathbb EX_n(p)}\longrightarrow v(p)/2.
\end{equation}

We next justify extension uniformly in $n$. For real $g,h$ put
$a=g-h$ and $g_u=h+ua$, $0\le u\le1$. Differentiation of the
finite determinant and then integration along its tilt path give
\[
 \frac{d}{du}G_n(g_u)
  =\mathbb E_{g_u}X_n(a)-\mathbb E_0X_n(a)
  =\int_0^1\Cov_{rg_u}(X_n(a),X_n(g_u))\,dr.
\]
Tilted covariance Cauchy--Schwarz and the preceding ordered identity
show
\begin{equation}\label{Gcontinuity}
 \left|\frac{d}{du}G_n(g_u)\right|
 \le e^{\operatorname{osc}(g_u)}
       \sqrt{\Var_n(a)\Var_n(g_u)}.
\end{equation}
By \eqref{varianceholder}, $G_n$ are uniformly Lipschitz on bounded
real $C^\eta$ sets for $\eta>1/2$.

Choose $1/2<\eta<\gamma$. Extend $F$ evenly to the circle.
Its Fej\'er convolutions $F_m$ are even trigonometric polynomials,
converge uniformly, and have bounded $C^\gamma$ seminorms.
The elementary increment estimate
\[
 [H]_{C^\eta}\le C\|H\|_\infty^{1-\eta/\gamma}
                        [H]_{C^\gamma}^{\eta/\gamma}
\]
proves convergence in $C^\eta$. Subtract
$a_m+b_m\cos\theta$ with
$a_m=(F_m(0)+F_m(\pi))/2$ and
$b_m=(F_m(0)-F_m(\pi))/2$. The corrected functions $Q_m$ vanish
at both endpoints, are polynomials in $\cos\theta$, and converge
to $F$ in $C^\eta$, since $a_m,b_m\to0$.

The circle $H^{1/2}$ difference-quotient identity, first verified by
Fourier orthogonality for trigonometric polynomials, identifies
$\sum |k||\widehat H_k|^2$ with a fixed positive multiple of
\[
 \iint\frac{|H(\theta)-H(\phi)|^2}
                   {\sin^2((\theta-\phi)/2)}\,d\theta\,d\phi.
\]
It is bounded by $C_\eta[H]_{C^\eta}^2$. Thus $v(F-Q_m)\to0$ and
$v(Q_m)\to v(F)$ by weighted Cauchy--Schwarz. Equations
\eqref{polynomialG} and \eqref{Gcontinuity}, first with $n\to\infty$
and then $m\to\infty$, give $G_n(g)\to v(F)/2$.
Finally the earlier Jacobi mean theorem applies because
$F=O(d^\gamma)$ at the endpoints. It gives
$\mathbb EX_n(g)=(n+1)\mu(F)+o(1)$, proving
\eqref{scalarconstant}.
\end{proof}

Apply this proposition to
\[
 F_\delta(\theta)=\Log(1+\delta w(\theta)),\qquad
 F_{\delta,k}=\frac1\pi\int_0^\pi F_\delta(\theta)\cos(k\theta)\,d\theta.
\]
For real $\delta>-3$ it is real, endpoint-vanishing and $C^{2/3}$,
and the exact pressure integral from \cite[Section 9]{r124} is
$\mu(F_\delta)=2f(3+\delta)$. Therefore
\begin{equation}\label{scalarreallimit}
 \log T_n(\delta)-2nf(3+\delta)
 \longrightarrow 2f(3+\delta)+\tfrac12\sum_{k\ge1}kF_{\delta,k}^2.
\end{equation}
The earlier complex scalar bound makes the left side locally bounded
and holomorphic on $\Omega$. Vitali's theorem extends this convergence
locally uniformly over $\Omega$. The displayed series is its analytic
expression there, with squares, not absolute squares. Indeed for
$1/2<\eta<2/3$ the map $\delta\mapsto F_\delta$ is holomorphic
as a $C^\eta$-valued map, by the uniform lower bound for
$|1+\delta u|$, $0\le u\le1/3$, on compact parameter sets.
Its image on a compact set is compact in the weighted Fourier
$\ell^2$ space. Finite-coordinate projections then have uniformly
vanishing tails on this image. The bilinear Cauchy--Schwarz inequality
proves local uniform convergence and holomorphy of the square series.
The real identity and analytic uniqueness prove the claimed expression.

\subsection{The explicit limit and branch normalization}
Combine \eqref{inheritedfactors}, \eqref{mixedlimit}, and
\eqref{scalarreallimit}. Every constituent logarithm is normalized at
$t=3$, and their locally uniform limits give
\begin{equation}\label{Eexplicit}
 \boxed{\ E(t)=\Log\mathcal A(t)+\tfrac12\Log\Xi(t)
       -\tfrac12\Log\Psi(t-3)+f(t)
       +\tfrac14\sum_{k\ge1}kF_{t-3,k}^2.\ }
\end{equation}
This proves $E_n\to E$ locally uniformly on $\Omega_t$.
At $t=3$ every term is zero. For positive $t$ the three determinant
factors are positive, so $E(t)$ is finite and real. Notice both the
minus sign of the mixed factor and the scalar coefficient $1/4$:
one first solves $S_n/Q_n^2\to\Psi$ for $\log Q_n$, and then
halves the scalar variance term. These fix the amplitude normalization.
As a separate normalization check, for the polynomial $p(x)=x$ the
only positive cosine coefficient is $1/2$, so $v(p)=1/4$ and the
centered logarithmic determinant tends to $1/8$. Directly the finite
variance is the square of the single Jacobi recurrence coefficient
crossing the degree-$n$ boundary, which tends to $1/4$.
